% Auto-generated by paper/build_survey.py -- do not edit by hand.
{\footnotesize
\setlength{\tabcolsep}{3pt}
\renewcommand{\arraystretch}{1.15}
\begin{longtable}{@{}p{2.1cm} >{\raggedright\arraybackslash}p{1.9cm} >{\raggedright\arraybackslash}p{1.5cm} >{\raggedright\arraybackslash}p{1.3cm} >{\raggedright\arraybackslash}p{1.5cm} >{\raggedright\arraybackslash}p{2.4cm} >{\raggedright\arraybackslash}p{3.0cm}@{}}
\caption{Processing choices of the 103 surveyed \dvv\ studies (the literature survey underpinning this paper). Every study is cited; \texttt{n/r} = not reported in the source. Frequency band, coda window, estimator, reference/stack scheme and uncertainty treatment are the choices Section~\ref{sec:results} shows to control the result.}\label{tab:survey}\\
\toprule
\textbf{Study} & \textbf{Site / region} & \textbf{Freq (Hz)} & \textbf{Coda (s)} & \textbf{Method} & \textbf{Reference / stack} & \textbf{Uncertainty treatment} \\
\midrule
\endfirsthead
\multicolumn{7}{@{}l}{\footnotesize\itshape Table~\ref{tab:survey} continued}\\
\toprule
\textbf{Study} & \textbf{Site / region} & \textbf{Freq (Hz)} & \textbf{Coda (s)} & \textbf{Method} & \textbf{Reference / stack} & \textbf{Uncertainty treatment} \\
\midrule
\endhead
\midrule\multicolumn{7}{r@{}}{\footnotesize\itshape continued on next page}\\
\endfoot
\bottomrule
\endlastfoot
\citet{Brenguier2008} & Piton de la Fournaise, La Réunion & 0.5–2 & n/r & Stretching & Daily stacks referenced to long-term stack & Compared against \textasciitilde{}0.05\% measurement precision / noise floor \\
\citet{Duputel2009} & Piton de la Fournaise, La Réunion & 0.1–1 & n/r & Other & Quasi-real-time stacks vs reference & n/r \\
\citet{Obermann2013} & Piton de la Fournaise, La Réunion & >0.5 & increasing lapse-time coda windows & Stretching & Daily stacks, referenced; coda at increasing lapse times & n/r \\
\citet{HotovecEllis2014} & Mount St. Helens, WA, USA & >5 (high-frequency repeating earthquakes) & n/r & Other & Stack across several hundred repeating-earthquake families (multiplets) & n/r \\
\citet{HotovecEllis2015} & Mount St. Helens, WA, USA & n/r & n/r & Other & Combines repeating-earthquake CWI with ambient-noise interferometry & n/r \\
\citet{Rivet2015} & Piton de la Fournaise, La Réunion & n/r & n/r & Stretching & Daily stacks vs reference; network-averaged dv/v & dv/v corrected for rainfall pore-pressure to improve detectability of small drops \\
\citet{DePlaen2016} & Piton de la Fournaise, La Réunion (method also re Kawah Ijen) & 0.1–1.0 (drop visible); best SNR 0.5–1 and 1–2 & n/r & Stretching & n/r & n/r \\
\citet{Donaldson2017} & Kīlauea summit, Hawaii, USA & 0.33–1.0 (lava-lake tremor band) & up to \textasciitilde{}120 (min lag set by inter-station distance to tremor source) & MWCS & Daily NCFs stacked over 3-day moving window & Reject pairs when MWCS regression error too large; accept if coherence high (>\textasciitilde{}0.65) \\
\citet{Takano2017} & Izu-Oshima, Japan & 0.5–1; 1–2; 2–4 (depth discrimination) & n/r & n/r & Cross-correlations 2012–2015, multiple bands & n/r \\
\citet{Lesage2018} & Volcán de Colima, Mexico & 0.125–2 & n/r & Stretching & Daily cross-correlations; 2013 stack as reference & Pair averaging → <0.05\% noise floor; clock-sync errors identified and corrected \\
\citet{DePlaen2019} & Mt. Etna, Italy (2013–2014) & 1–2 & 5–35 (causal + acausal lags) & MWCS & Daily autocorrelations, 2-day linear stack & Median errors \textasciitilde{}0.03\%/station; reject if dt error >0.1 s and coherence <0.6 \\
\citet{Donaldson2019} & Northern Volcanic Zone (Askja/Bárðarbunga), Iceland & multiple bands for depth discrimination (n/r exact) & n/r & n/r & n/r & n/r \\
\citet{Olivier2019} & Kīlauea, Hawaii, USA (2018 eruption) & 0.1–0.2 (secondary microseism); 0.3–1 (volcanic tremor) & n/r & Stretching & n/r (passive image interferometry, daily) & Resolves <0.1\%; cross-checked vs GPS/tilt \\
\citet{Yates2019} & Whakaari / White Island, New Zealand & n/r & n/r & n/r & n/r & n/r \\
\citet{Feng2020} & Kīlauea, Hawaii, USA & 0.3–0.6; 0.6–0.9; 0.9–2.0 & n/r & n/r & Jan 2017–Jun 2018 NCFs & n/r \\
\citet{HotovecEllis2022} & Kīlauea, Hawaii, USA (2018 collapse) & n/r & n/r & Other & Coda Wave Interferometry with Repeating Earthquakes (CWIRE), per-cluster pairs & n/r \\
\citet{Kopfli2024} & Mount St. Helens, WA, USA & n/r & n/r & Other & Long-term NCF stacks; many dv/v series on a common grid & n/r \\
\citet{Yates2024} & Mt. Ruapehu, New Zealand (2005–2009) & n/r & n/r & Wavelet & n/r & n/r \\
\citet{Yukutake2025} & Izu-Oshima, Japan (2003–2020) & 0.5–2.0 & n/r & n/r & Long-term interferometry stacks, 2003–2020 & n/r \\
\citet{Schaff2004} & 1989 Mw 6.9 Loma Prieta aftershock zone, California (also Parkfield repeaters) & n/r & n/r & Other & Repeating-earthquake doublets; delays vs pre-mainshock repeaters & High-precision relative delay via waveform cross-correlation of near-identical repeaters \\
\citet{Pacheco2005} & Theory / acoustic diffusion model & n/r & n/r & Other & n/r & n/r (analytical kernel) \\
\citet{Rubinstein2005} & 2004 M6 Parkfield earthquake, San Andreas fault, California & n/r & n/r & Other & Repeating-earthquake travel-time delays before vs after mainshock & Precise relative delay from near-identical repeating-earthquake waveforms \\
\citet{Wegler2007} & Mid-Niigata (Chuetsu) earthquake source region, Japan; F-Net station KZK (\textasciitilde{}24 km) & 2–50 (2 Hz high-pass; 100 Hz data) & 5–14 & Stretching & Daily autocorrelations vs \textasciitilde{}two-week pre-event reference; grid search (10000 trials, dv/v… & ±0.1\% fluctuation = inversion accuracy \\
\citet{Brenguier2008b} & Parkfield, San Andreas fault, California, USA & 0.1–0.9 & n/r & Stretching & 1550-day reference vs overlapping 5-day moving windows; 30-day time resolution & Averaged over 78 receiver pairs plus 5-day segments \\
\citet{Hadziioannou2009} & Laboratory scattering gel (ultrasonic) & 2.5e6 (2.5 MHz lab) & n/r & Other & Noise-correlation stacking & Demonstrates stability when background noise-source structure is stable \\
\citet{Sawazaki2009} & Western-Tottori region, Japan; KiK-net borehole (surface + 100 m); 2000 Mw 6.6 mainshock & n/r & n/r & Other & Coda deconvolution of surface vs 100-m borehole records & n/r \\
\citet{Wegler2009} & 2004 Mw 6.6 Mid-Niigata earthquake source region, Japan; 6 Hi-net/F-Net stations <25 km & n/r & n/r & Stretching & Daily auto-/cross-correlations inverted daily for 2 months pre/post; pre-event reference & Daily inversion fluctuations; consistency across 6 stations \\
\citet{Chen2010} & Sichuan, China — 2008 Mw 7.9 Wenchuan earthquake fault zone & n/r & n/r & Other & Noise cross-correlation stacking; sub-array comparison (156 broadband stations) & Pre-earthquake fluctuation (\textasciitilde{}0.02\%) used as baseline/noise reference \\
\citet{Nakata2011} & NE Japan (Honshu), KiK-net; 2011 Tohoku-Oki earthquake & n/r & n/r & Other & >300 earthquakes (Jan 1–May 26 2011); borehole-to-surface deconvolution & n/r \\
\citet{Rivet2011} & Guerrero, Mexico subduction zone (2006 M7.5 slow slip event) & 0.06–0.14 (periods 7–17 s) & n/r & MWCS & n/r & n/r \\
\citet{Hobiger2012} & 2008 Iwate-Miyagi Nairiku earthquake (Mw 6.9), NE Japan & 0.125–1.0 (three sub-bands) & n/r & Stretching & Daily cross-correlations vs long-term reference Green's function & Curves fit by model (coseismic drop + log postseismic recovery + seasonal) \\
\citet{Minato2012} & Southern Tohoku (Fukushima/Ibaraki), Japan; 58 Hi-net stations; 2011 Tohoku-Oki & n/r & n/r & Stretching & n/r & n/r \\
\citet{Obermann2013} & Numerical 2-D elastic media (lunar-data illustration) & central 20 (bandwidth \textasciitilde{}12) & 1.8–6.6 (primary lapse \textasciitilde{}3.6 s) & Stretching & Averaged over 10 random medium realizations per configuration & S.d. over 10 realizations; mean-free-path uncertainty \textasciitilde{}2–15\% \\
\citet{Brenguier2014} & Volcanic/geothermal regions of Honshu, Japan; 2011 Mw 9.0 Tohoku-Oki & 0.1–0.9 & n/r & MWCS & Daily cross-correlations; linear inversion/regularization for continuous time series & Stress sensitivity \textasciitilde{}0.001 MPa\textasciicircum{}-1 inferred; detailed dv/v uncertainty n/r \\
\citet{Liu2014} & Epicentral region of 2008 Mw 7.9 Wenchuan earthquake, China & 0.25–0.5 (period 2–4 s most sensitive) & n/r & Stretching & Noise cross-correlation stacking (Aug 2004 – Sep 2011) & Pre-event period as baseline (no obvious precursor) \\
\citet{Taira2015} & South Napa, California (2014 Mw 6.0 South Napa earthquake) & 0.08–2.0 & n/r & MWCS & MSNoise cross-correlations; decimated to 20 Hz & n/r \\
\citet{Gassenmeier2016} & Station PATCX, Atacama Desert, Northern Chile (IPOC) & 4–6 primary (also 1–3, 7–10) & 10–15 primary (also 5–10, 15–20) & Stretching & Daily autocorrelations; two-step iterative reference & Error from correlation-coefficient threshold ±1 s.d.; healing-fit errors (amplitude \textasciitilde{}21\%… \\
\citet{Hillers2019} & San Jacinto fault zone, California (after 2010 M7.2 El Mayor-Cucapah, M5.4 Collins Valley) & n/r & n/r & Stretching & Average-waveform reference for stretching; MWCS also applied & Cross-check between time-domain stretching and frequency-domain MWCS \\
\citet{Wang2019} & NE Honshu, Japan (2011 Mw 9.0 Tohoku-Oki earthquake) & 0.02–1.0 (period 1–50 s; short-period 1–7 s + tiltmeters 8–50 s) & n/r & Stretching & Monthly velocity-change estimates & n/r \\
\citet{Mao2020} & Salton Sea Geothermal Field, California (2009–2011) & frequency-dependent across band (time-frequency domain) & full time-frequency domain & Wavelet & Noise cross-correlation stacking & Wavelet-coherence thresholding and amplitude-based weighting \\
\citet{Poli2020} & L'Aquila region, central Italy (2009 Mw 6.3 L'Aquila earthquake) & n/r & n/r & n/r & n/r & n/r \\
\citet{Boschelli2021} & Ridgecrest fault zone, California (2019 Mw 7.1 Ridgecrest earthquake) & n/r & n/r & Other & Daily autocorrelation functions vs mean waveform & n/r \\
\citet{Lu2021} & Ridgecrest, California (2019 Mw 7.1 Ridgecrest earthquake) & 8.0–12.0 & 3 (moving windows, 1.5 s step) & Other & 10-min stacks; adaptive Gaussian smoothing (20 min growing to 24 hr over 3 days) & Standard deviations tracked; stabilized via smoothing \\
\citet{Sheng2022} & San Jacinto fault, Anza seismic gap, Southern California & >1.0 & n/r & Wavelet & Weekly stacked correlations with \textasciitilde{}2-month smoothing & Inverse travel-time uncertainties used as weights in misfit \\
\citet{Mainsant2012} & Pont-Bourquin landslide, Swiss Alps (Switzerland) & \textasciitilde{}4-10 (band of measured surface-wave dv/v; spectral analysis up to \textasciitilde{}12) & n/r (short inter-sensor distance, coda of daily correlograms) & Stretching & Daily noise correlograms (cross-correlation functions) & Daily correlation-coefficient / waveform-coherence checks; qualitative \\
\citet{Voisin2016} & Avignonet / Mas d'Avignonet landslide, Trièves (French Alps), France & \textasciitilde{}6-8 (annual dv/v pattern best resolved) & n/r & Stretching & Daily/seasonal stacks vs reference & Correlation-coefficient weighting; comparison with piezometer \\
\citet{Harba2017} & Just-Tęgoborze landslide, Carpathian flysch, southern Poland & High-frequency seismic noise (\textasciitilde{}5-20+, dispersion-curve range) & n/r (interferometric Green's-function retrieval) & Other & Cross-correlation / interferometric stacks; dispersion inversion (neighbourhood algorithm) & Dispersion-inversion misfit; comparison with MASW \\
\citet{Bertello2018} & Montaguto earthflow, southern Italy & n/r (surface-wave / MASW band, \textasciitilde{}few-20) & n/r & Other & Time-lapse interferometry + time-lapse active MASW & Comparison passive vs active MASW; qualitative \\
\citet{Bievre2018} & Pont-Bourquin landslide, Swiss Alps (4.5-yr record) & \textasciitilde{}5-15 (surface-wave band; reported sensitivity to shallow layer <=2 m) & n/r & Stretching & Daily correlations, multi-year reference & Correlation coefficient; seasonal-cycle modeling to separate from precursors \\
\citet{Colombero2018} & Madonna del Sasso cliff, NW Italian Alps (Orta Lake), Italy & 2-20 (analysis); 2-4 strongest annual signal & 0.5-2 (and symmetric -2 to -0.5) & Stretching & Daily correlations, multi-year reference (2013-2016) & Correlation coefficient (\textasciitilde{}0.9 at low freq) as reliability metric \\
\citet{Bontemps2020} & Maca slow-moving Andean landslide, southern Peru (3-yr dataset) & \textasciitilde{}1-10 (landslide-scale) & n/r & Stretching & Daily stacks vs reference & Confidence bands; gap-flagged periods; correlation with rainfall and seismicity \\
\citet{Fiolleau2020} & Harmalière landslide, French Western Alps (collapsed Nov 2016) & 1-12 (monitoring); dv/v in 2-4 and 8-10 bands; block resonance \textasciitilde{}9-16 & 0.05-1.5 & Stretching & Hourly records, daily averaging vs reference & Correlation-coefficient degradation tracked across bands; multi-parameter precursor timing \\
\citet{LeBreton2021} & Global review (9 landslides monitored to date) & \textasciitilde{}1-20 across studies (landslides typically higher than tectonic) & Sub-second to a few s (short inter-station distances) & Stretching & Daily stacks vs long-term reference & Reviews spurious fluctuations (noise-source, site-response, inter-sensor distance change… \\
\citet{Fan2023} & Rock slope, southwest China & High-frequency (\textasciitilde{}5-20) & Sub-second to a few s & Stretching & Short-window stacks for rapid detection & Decorrelation/CC checks; rapid-detection confidence discussion \\
\citet{Liu2025} & Xishan Village landslide, Sichuan Province, China & 5-25 (raw 1-30); 15 Hz top \textasciitilde{}15 m, 5 Hz 15-40 m & 0.4-2.0 & Stretching & 20-min interval CCFs vs reference (SNR>4 wave-packet selection) & 95\% confidence interval (\textasciitilde{}0.05\% uncertainty at 5 Hz) \\
\citet{Whiteley2026} & Hollin Hill Landslide Observatory, North Yorkshire, UK (2-yr) & n/r (near-surface surface-wave / MASW band, \textasciitilde{}5-30) & n/r & Other & Time-lapse interferometry + MASW between recording periods & Statistical distribution of surface-wave velocity measurements over time \\
\citet{deWit2026} & Open-pit mine slope, Australia & n/r (high-frequency near-surface) & n/r & Stretching & Time-lapse stacks vs reference & Decorrelation/CC; comparison with radar surface deformation \\
\citet{SensSchonfelder2006} & Merapi Volcano, Indonesia & 0.5-10 (analysis >0.5) & 2-4 and 6-8 & Stretching & Daily autocorr vs yearly reference & Correlation-coefficient maximization; \textasciitilde{}0.1\% accuracy quoted \\
\citet{Tsai2011} & Theoretical model, applied to southern California & n/r & n/r & Other & n/r & Analytic sensitivity / parameter study \\
\citet{Hillers2014} & TCDP borehole array, Chelungpu fault, Taiwan & n/r & n/r & n/r & n/r & n/r \\
\citet{Lecocq2017} & Gräfenberg Array (GRA1-4), SE Germany (karst limestone aquifer) & 0.1-0.8 (\textasciitilde{}7 s secondary microseism) & 20-100 lag; 12 s window, 2 s step (MWCS) & MWCS & Daily CCFs, 31-day rolling window & Coherency-weighted least squares; MWCS errors; jackknife stability \\
\citet{Nimiya2017} & Kyushu Island, Japan (2016 Kumamoto region) & 0.1-0.9 & \textasciitilde{}100 (MWCS coda-onset to 60 s lapse) & Stretching & Daily CCFs from 30-min segments; 1-yr reference, 30-day moving stack & Stretching vs MWCS cross-check; inter-station distance <40 km \\
\citet{Clements2018} & San Gabriel Valley, California, USA & 2-4 (0.2-0.5 band discarded as too noisy) & n/r (\textasciitilde{}1-3 consistent with 2-4 Hz) & Stretching & Daily stacks vs reference & Stretching correlation-coefficient; comparison to GPS/well data \\
\citet{Kim2019} & Gulf Coast Aquifer, southern Texas (IU.HKT near Houston), USA & n/r & n/r & Other & Multi-year (\textasciitilde{}20 yr) record; yearly-scale comparisons & n/r \\
\citet{Andajani2020} & Chugoku \& Shikoku, SW Japan (Hi-net) & \textasciitilde{}0.1-0.9 (low-frequency emphasized) & n/r & Stretching & n/r & Band-pass separation of rainfall vs sea-level/atmospheric forcing; lag analysis \\
\citet{GaubertBastide2022} & Crépieux-Charmy water-production field, Lyon, France & \textasciitilde{}2-4 (little energy >10 Hz) & n/r & Stretching & Hourly correlations over 19 days, two fill/drain cycles & Correlation-coefficient QC; dense-array spatial consistency \\
\citet{Illien2022} & Nepal Himalaya (2015 Mw7.8 Gorkha aftermath) & n/r & n/r & Stretching & \textasciitilde{}3-yr continuous time series & Coupled seismo-hydrological model \\
\citet{Mao2022} & Coastal Los Angeles basins, California, USA & 0.2-2.0 (sub-bands 0.2-0.8, 0.4-1.6, 0.5-2.0) & Early coda \textasciitilde{}15-50 & Other & Daily CCFs stacked over 20 days, 5-day step; pairs <50 km & Bayesian inversion; covariance from CC coefficients + smoothing matrix \\
\citet{Clements2023} & Statewide California, USA (1999-2021) & 2-4 & n/r (\textasciitilde{}1-3 consistent with 2-4 Hz) & Stretching & Daily stacks across \textasciitilde{}700 stations vs reference & Stretching correlation-coefficient; drained-poroelastic rainwater-diffusion model \\
\citet{Delouche2023} & Greece (aquifer monitoring sites) & n/r & n/r & Stretching & n/r & Linear precipitation model; correct dv/v for hydrologic contribution before tectonic int… \\
\citet{Ermert2023} & Mexico City basin / Valley of Mexico & 0.5-1, 1-2, 2-4, 4-8 (octave bands) & n/r & Stretching & Clustered (GMM) autocorrelation stacks & MCMC inversion; 16th-84th percentile; CC>0.6 retained \\
\citet{Fokker2023} & Groningen, The Netherlands & n/r (frequency-dependent surface-wave kernels) & n/r & Other & n/r & Inversion of frequency-dependent dv/v to depth-resolved pore pressure \\
\citet{Zhang2023} & Central Oklahoma, USA & n/r & n/r & n/r & Continuous 2013-2022 record & n/r \\
\citet{Mao2025} & Greater Los Angeles, California, USA & \textasciitilde{}0.2-2.0 (multi-band, as Mao 2022) & Early coda \textasciitilde{}15-50 & Other & Daily CCFs stacked, multi-day step; two-decade record & Bayesian depth inversion (as Mao 2022) \\
\citet{Mordret2016} & Southwest Greenland ice sheet & n/r & n/r & Stretching & Daily cross-correlations over \textasciitilde{}2-year record & n/r \\
\citet{James2019} & Interior Alaska (Poker Flat / Fairbanks area), USA & 3-30 (multiple sub-bands) & n/r & Stretching & Daily cross-correlations stacked; reference relative to Jan 2014 & Correlation-coefficient quality control; Bayesian tomography for spatial maps \\
\citet{Guillemot2020} & Gugla rock glacier, Valais, Switzerland & 1.5-14 (bands 4-8, 7-11, 10-14) & 0.3-0.8 (causal and acausal) & Stretching & Daily correlograms from hourly raw data & Weaver et al. (2011) error formula \\
\citet{Guillemot2021} & Laurichard (French Alps) and Gugla (Swiss Alps) rock glaciers & 1-50 (resonance peaks >10) & n/r (spectral/PSD method, not coda) & Other & PSD-based; finite-element modal modeling & Modal modeling comparison \\
\citet{Lindner2021} & Mt. Zugspitze, German/Austrian Alps & n/r (high-frequency local band) & n/r & Stretching & 15-year continuous record; daily/seasonal stacks & Comparison with meteorological/borehole records; correlation coefficient \\
\citet{Luo2023} & Greenland Ice Sheet (GLISN stations) & n/r & n/r & Stretching & Daily autocorrelations stacked & Lag-time correlation with GRACE/mass data \\
\citet{Gassenmeier2015} & Ketzin CO2 storage site (CO2SINK), Brandenburg, Germany & 1.0-3.5 (detail 1.5-3) & Windows after 300 m/s arrival (lapse-time dependent) & Stretching & 1-hour segments cross-correlated, stacked to daily & Weaver et al. (2011) errors; confidence intervals \\
\citet{Hillers2015} & Basel deep geothermal (EGS) reservoir, Switzerland & n/r & Coda of noise correlations & Stretching & Daily noise correlations around 2006 stimulation & Sensitivity-kernel imaging of velocity-change location \\
\citet{Obermann2015} & St. Gallen geothermal site, Switzerland & n/r (sub-Hz to few Hz local) & Coda of noise cross-correlations (lapse-time dependent) & Stretching & Daily cross-correlations; reference stacks & Sensitivity-kernel-based localization; noise-source checks \\
\citet{Czarny2016} & Underground coal mine, Upper Silesia, Poland & 0.6-1.2 & n/r & Stretching & Continuous cross-correlations over \textasciitilde{}42 days & n/r \\
\citet{Olivier2017} & Active tailings storage facility (mine), South Africa & \textasciitilde{}5-15 & n/r & Stretching & Daily cross-correlations across array & n/r \\
\citet{Taira2018} & Salton Sea Geothermal Field, California, USA & 0.25-8.0 (six bands) & 20 (coda windows -40 to -20 s and 20 to 40 s) & MWCS & Daily NCFs from 30-min segments; 5-day stacks & Two-sigma std; correlation-coefficient threshold >0.85 \\
\citet{Kristjansdottir2019} & Hellisheidi geothermal field, SW Iceland & n/r & n/r & Stretching & Daily correlations (MSNoise workflow) & Seasonal and noise-source effects discussed as caveats \\
\citet{Snieder2002} & theory + ultrasonic lab & n/r & later coda (multiply scattered arrivals) & Other & n/r & variance of time shift from cross-correlation; statistical theory of multiply scattered… \\
\citet{Bensen2007} & USArray test data & whitening band stated (e.g. 5 Hz-100 s periods) & n/r (surface-wave window) & n/r & daily cross-correlation + temporal stacking & SNR-based quality control; error analysis on dispersion measurements \\
\citet{Clarke2011} & synthetic + real noise tests & n/r & moving windows in coda & MWCS & current stack vs reference, weighted linear regression dt vs t & formal MWCS errors: coherence-weighted phase-delay errors per window; error inversely pr… \\
\citet{Hadziioannou2011} & urban noise autocorrelations & n/r & coda of autocorrelations & Stretching & waveform clustering to enhance daily stacks & clustering reduces noise-source variability; resolution improved from 30 days to 1 day \\
\citet{Weaver2011} & theory & n/r & coda window (analytical) & Stretching & reference vs current correlation & analytic RMS error of stretching estimate as function of coherence/decorrelation, bandwi… \\
\citet{Zhan2013} & synthetic + field examples & sensitive to band choice & n/r & Stretching & reference vs current & identifies systematic (non-random) bias from non-stationary noise spectrum; recommends c… \\
\citet{Lecocq2014} & Piton de la Fournaise (validation) & user-defined (e.g. 0.1-1 Hz) & user-defined moving windows & MWCS & daily CCF + reference stack & implements Clarke et al. 2011 MWCS errors; reproducible parameter database \\
\citet{Mikesell2015} & numerical CWI examples & n/r & full-trace local time shifts & DTW & n/a & compares accuracy of WCC vs stretching vs DTW; DTW more robust at low SNR \\
\citet{Stehly2015} & Wenchuan Mw 7.9 region (test dataset) & n/r & coda of daily correlations & Stretching & daily correlations curvelet-filtered before stacking & curvelet denoising reduces dv/v uncertainty, enabling 1-day resolution \\
\citet{Daskalakis2016} & synthetic + real noise & n/r & coda window & Stretching & normalized cross-correlation functions & shows seasonal noise-source fluctuations create apparent dv/v; adequate CCF normalizatio… \\
\citet{Obermann2016} & 3-D wavefield simulations & frequency-dependent & lapse-time dependent & Other & n/a & n/r \\
\citet{Wang2017} & Japan / Hi-net (test dataset) & \textasciitilde{}0.1-0.9 Hz (multi-band) & coda windows & MWCS & compares stacking schemes; long reference stack for seasonal trends & demonstrates reference-stack and stacking-length choice control the apparent long-term t… \\
\citet{Obermann2019} & review & scale-dependent & lapse-time/depth dependent & Other & reference and moving-stack strategies reviewed & synthesizes error sources: noise-source variability, processing choices, sensitivity ker… \\
\citet{Jiang2020} & benchmark vs MSNoise & user-defined & user-defined & Other & linear/PWS/robust stacking; reference + moving stacks & offers multiple dv/v estimators so users can cross-check; parallel/HDF5 reproducible pip… \\
\citet{Wang2020} & review & application-dependent & application-dependent & Other & reference and stacking strategies reviewed & reviews uncertainty/bias sources and method trade-offs \\
\citet{Yuan2021} & 2-D heterogeneous half-space simulations & multiple bands tested & coda windows tested & Other & n/a & quantifies accuracy/robustness of 7 estimators vs SNR, perturbation size, heterogeneity;… \\
\end{longtable}
}
